\subsection{Charge, flux and filtered arrival on the interacting source}

The central hypercharge current supplies an exact field attachment on the
eight-cell source. Orient each stored link from its left to its right
endpoint and set $b_e=L^Y_{e,0}$. The complete 119-state link satisfies
$L^Y_{e,1}=-b_e$. With $(Df)_e=f_{\mathrm{right}}-f_{\mathrm{left}}$,
the physical Gauss constraints give
\[
 D^*b=\rho,\qquad
 \rho_x=\rho_x^{\mathrm{matter}}+\rho_x^{\mathrm{scalar}}
       +\delta_{x0}Y_{\mathrm{ref}}.
\]
On the invariant one-$e^c$-reference sector populated by the native removal
pump, each of the six elementary $e^c$ arrival routes has
$Q=Y=-1$ and the retained reference has $Y=+1$. Summing the actual
endpoint currents over any region gives its total enclosed charge,
including that reference. All $6\cdot256$ regional identities and all
$24{,}840$ native factor-column Gauss comparisons are checked. The full
flux spectra include every link state. The noncentral weak component of
$Q=T_3+Y$ does not admit the same endpoint cancellation on an arbitrary
link; the stated central-current identity does not omit its gauge charge.

An energy-resolved experiment uses the actual full-$H$ bin projector $E$
of radius $1/1024$ about the visible $e^c$ pole and an arrival projector
$F_e$ refined by its local atomic pole. If $a$ is the atomic pole energy,
then $H_{\rm at}F_e=(a+1)F_e$, where
$H_{\rm at}=H_0+4N_{\rm ref}$, and $F_eTF_e=0$. The latter follows from the
singlet compression of the other matter endpoint and holds with all link
states retained. The resolvent identity controls $F_eE$ without assuming
that the complete charged band is isolated. Together with the exact
native pole residue, it yields original-input probabilities between
$1.04092024\,10^{-23}$ and $1.68685307\,10^{-23}$ for each tagged
arrival at $t=0$. Their change at $t=2^{-24}$ is at most
$2.24579\,10^{-34}$. These records measure virtual dressing present in
the filtered state; they do not establish propagation from the origin.
All seventeen outcomes, including preparation, pump, bin, untagged and
nonarrival outcomes, are retained. The local tag can change the energy
expectation and is an active intervention.

This experiment records the atomic ground and blank-link projector $P_0$
before pumping and filtering, so its successful initial vector is
$EBP_0\Psi$. The finite excursion above instead measures $S_0$ after
energy selection, giving $S_0EB\Psi$ and exactly zero initial arrival.
Its positive later arrivals and quantified energy disturbance concern
that complete measurement sequence. The virtual initial tail and the
position-prepared excursion have distinct preparations and are not
interchangeable transport laws.

The reference also has a quantitative binding effect. Let $\Pi_y$ require
one unscreened charge at $y$ and gauge-singlet matter at every other cell,
retaining all links. Exact Gauss strings and local endpoint embeddings
give the constrained energy differences
\[
 \mathcal E_L(y)-\mathcal E_L(0)=d(0,y),\qquad
 d(0,y)-w\leq\mathcal E_{e^c}(y)-\mathcal E_{e^c}(0)
          \leq d(0,y)+w,
 \quad w\simeq0.591071386691.
\]
The additive reference cost cancels from these differences. The compressed
subspaces need not be invariant under $H$, and the complete complement
contains screening clouds. Thus this is a result about the specified
single-unscreened-charge family, with no assertion of unrestricted
confinement or free-electron identification.

\subsection{Native free beats with complete coherent heralding}

Let $P=P_0\otimes|0_{\rm ref}\rangle\langle0_{\rm ref}|$ restrict the
atomic source ground projector to the reference vacuum. The local pole
residue defines a partial isometry $A_r=S_rB_rP/\sqrt{w_r}$ with
$A_r^*A_r=P$ and $A_r^2=0$.
One declared coherent preparation makes an equal superposition of its
two source subspaces. Thereafter the unchanged full Hamiltonian evolves
the fixed quadratures $X=A_r+A_r^*$ and $Y=i(A_r-A_r^*)$. In the
atomic comparison their expectations are $\cos\Omega t$ and
$-\sin\Omega t$. For the ground-to-$e^c$ transition,
\[
 \Omega=4+\Delta_{e^c}\simeq4.691656680938395.
\]
The static reference energy is part of the measured beat. A common
herald onto the union of its two actual full-$H$ bins preserves
coherence. If $\eta$ bounds the projector error, the waveform error
is at most $2\eta+2v|t|$, uniformly over localized spins and untouched
references. At $t=10^6$ it is below $0.041953$.

The distinct reference families also support
$A_{e^c}A_L^*$, whose arms each contain one reference particle.
Its frequency is
$\Delta_{e^c}-\Delta_L\simeq0.594768183046114$, with exact
cancellation of the common reference cost. Resolving the two herald
arms separately destroys the relevant coherence, as does
reference-family dephasing. Seven complete records retain the
preparation and herald failures and the quadrature outcomes. The
initial coherent preparation and its work are explicit apparatus
inputs; these free beats supply neither passive self-starting nor a
physical Compton-clock identification.

\subsection{Coherent compilation of the lepton rest Hamiltonian}

The full native pairing admits a finite-group reduction that preserves
the entire physical Fock carrier. If $N_A$ counts ancillary fermions,
put $R=\exp(i\pi N_A/2)$ and
\[
 \overline V=\frac14\sum_{r=0}^3R^{-r}VR^r.
\]
The average removes every monomial creating or annihilating one or two
ancillary fermions. For the vacuum-ancilla embedding $J$ it obeys
$\overline VJ=JH_Y$, where $H_Y$ is the analyzer-derived mass pairing.
Two-phase parity averaging fails on ancillary pairs. Parity of the
complete charged-lepton leg changes the sign of $V$ and implements
negative interaction segments without discarding any Fock state.

Let $S_2$ be the symmetric seven-factor formula for the four conjugates
$R^{-r}VR^r/4$. With $p=(4-\sqrt[3]{4})^{-1}$, the composition
\[
 S_4(h)=S_2(ph)^2S_2((1-4p)h)S_2(ph)^2
\]
matches the exponential through degree four. Exact noncommuting-word
coefficients verify the cancellation. Since $\|V\|\leq6s$ on $s$
rest sites, a Taylor-tail bound and unitary telescoping give complete,
unhalved output trace error
\[
 \epsilon\leq\frac{(14s)^5|t|^5}{10n^4},\qquad n\geq14s|t|,
\]
uniformly over all inputs and retained quantum references. At $s=t=1$,
$n=15{,}229$ cycles, comprising $533{,}015$ interaction segments,
suffice for $\epsilon\leq10^{-12}$. The repeated-reset implementation
has a conservative budget of $108{,}000{,}000{,}000{,}000$ rounds at
that tolerance. The coherent construction uses one initial ancillary
vacuum and retains every final outcome. Its interaction time is
$(8p-1)|t|$, and its phase and parity controls are declared services.
The proof assumes zero kinetic Hamiltonian; it does not reverse a
noncommuting kinetic term.

The three coefficients, approximately $1:206.7703:3477.4728$, are
unchanged by this compilation. With the empirical electron anchor
$0.00051099895069\,\mathrm{GeV}$, the guarded numerical coefficient
intervals are
\[
 \begin{aligned}
 m_\mu&\in[0.105659408862295,0.105659420129515]\,\mathrm{GeV},\\
 m_\tau&\in[1.77698487649602,1.77698506513085]\,\mathrm{GeV}.
 \end{aligned}
\]
These intervals enclose the compiled coefficients under the declared
register and calibration. They are not physical uncertainty intervals;
the retrospective register choice, approximately $9.83$ parts-per-million
muon discrepancy, interacting-particle identity and physical clock
remain distinct scientific conditions.

\subsection{A neutral protected sector and its boundary dependence}

The actual Higgs-lepton composite
$J_s=\phi_0c_{s,11}-\phi_1c_{s,10}$ commutes with all twelve native
gauge currents. It is odd under fermion parity and changes raw lepton
number by one. Raw lepton number and baryon number commute with the
unchanged local Yukawa and hopping Hamiltonian. The full atomic gap
$g$, hopping bound $v=360\,2^{-34}$ and an admissible local neutral
trial give
\[
 g-v\leq E_{\mathrm{sector,min}}-E_{\mathrm{ground}}
       \leq\Delta_\nu+v,
 \qquad \Delta_\nu\simeq0.089555330476.
\]
The numerical enclosure is $[0.04117056754057018,
0.08955535143083537]$ in source units. The lowest-sector projector is
stationary and cannot decay to the source ground under the specified
lepton-number-preserving dynamics. This ideal spectral projector is defined
mathematically; a full interacting eigenvector or a primitive
preparation of that projector is not computed here.

The original antiperiodic translations square to fermion parity and
exclude zero momentum in this odd sector. A separately declared periodic
member reverses twelve seam signs. The second-order ground operator
and quartic remainder bounds are unchanged. Eight translated normalized neutral
trials are orthogonal, so their zero-momentum average has squared norm
$1/8$ and supplies the same positive threshold bound. Using the
existing electron energy calibration gives a conditional interval
$[3.0416993564\,10^{-5},6.6163881404\,10^{-5}]\,\mathrm{GeV}$.
This is a neutral finite-sector energy, with no inferred observed
neutrino identity, relativistic mass shell or cosmological abundance.

\subsection{Stationary intrinsic magnetic response}

An exact unitary decomposes the complete sixteen native Clifford rows
into four Dirac blocks while retaining every valley label. A declared
spatial construction adds a scalar confining well and Wilson regulator
to give a Hermitian operator $D_W(A)$ with positive mass lower bound
$m$. Its positive matter kernel is $K(A)=|D_W(A)|$. For a normalized
one-particle state and a mean-zero transverse field, use
\[
 \mathcal E(\psi,A)=\langle\psi,K(A)\psi\rangle
              +\frac{a^3}{2\kappa}\|CA\|^2.
\]
Compactness of the finite state sphere and field coercivity prove that
a minimum exists. The derivative includes both signed matter branches:
if $X=\delta K$, then
$KX+XK=D_W\delta D_W+\delta D_WD_W$.
The Peierls derivative, including the Wilson contribution, supplies
the actual bond current. The stationary solve obeys
$D^*J=0$ and $C^*CA=\kappa J/a^3$; its current is recomputed from
its own field.

With $m=8$, scalar-well strength $0.01$, regulator $1/2$ and coupling
$\kappa=0.1$, three fixed-extent refinements give bound-energy
$g$ values $1.47516960$, $1.68416493$ and $1.78875065$.
The convention is $g=2E_{\mathrm{bound}}\mu_z/(qS_z)$ with
$q=-1$; replacing the measured bound energy by the supplied $m$
gives different numbers. The computed unregulated central-difference
control has intrinsic current below $10^{-10}$ on its selected spatial
parity branch.
The nonzero stationary dipole, its Maxwell residual and its static
weak-field energy derivative are numerical results. The displayed
branch is not certified to be the global minimum, and three spacings
do not establish a continuum value of two. The construction includes
external mechanical confinement and a semiclassical transverse field;
it supplies no longitudinal electromagnetic completion or executed
time-dependent Larmor experiment.
Its classical connection and transverse Maxwell energy belong to this
distinct construction. They do not satisfy the independent quantum-link
conjugation hypotheses of the interacting charged source.

\subsection{Longitudinal field energy on the native charge carrier}

The same native Gauss identity permits the readout
\[
 G=(D^*D)^+,\qquad E_{\mathrm{long}}=DG\rho,
 \qquad \tfrac12\|E_{\mathrm{long}}\|^2=\tfrac12\rho^*G\rho.
\]
It is the minimum-norm field with divergence $\rho$. The complete
native flux retains its divergence-free part; this readout does not
project a quantum state onto a different carrier. Each of the six
unit-charge arrivals on the side-two cube has quadratic value $7/48$.
On the original source this is an observable, rather than its link
energy. Adding this positive energy to the unit-link penalty cannot
cancel the constrained linear binding cost.

A separate declared Hamiltonian replaces that penalty:
\[
 H'=H_{\mathrm{local}}+\lambda T+4N_{\mathrm{ref}}
                         +\frac\kappa2\rho^*G\rho.
\]
All matter, scalar, reference and link states are retained. The central
density commutes with the local gauge action, so all Gauss constraints
and the native continuity equation are preserved. Choose the native
$e^c$ local endpoint eigenstate at $y$, a local ground state at every
other cell, a fixed one-particle reference at zero and a normalized
physical flux path connecting the endpoints. This trial has exact expectation
\[
 (V-1)E_0+e_{e^c}+4+\kappa(G_{00}-G_{0y}),
\]
independent of its path length. The same local lower bound gives a
constrained energy enclosure of distance-independent width. The
computed Green interaction approaches $-\kappa/(4\pi r)$ in counting
distance; finite-box potential-difference and force amplitude ratios
at side 128 and radius eight are approximately $1.00407$ and $1.02842$.
These are static trial and field energies. At nonzero hopping the
strings need not be eigenstates. Unpenalized flux loops retain a
degeneracy, and the original ground preparation and clock do not
transfer to $H'$. The coupling and physical length normalization are
declared; these identities determine no physical fine-structure constant.

The static hypercharge obstruction also holds for this alternative
Hamiltonian. Its density acts on matter, scalar and reference factors,
so $[K_e,\rho_x]=0$ and the longitudinal energy commutes with every
$V(A)$. Consequently $H'(A)=V(A)H'(0)V(A)^\dagger$: replacing the
linear link penalty by the Coulomb energy does not produce static
hypercharge spectral splitting. This operator identity imports no ground
preparation, charged energy bin or transport law from $H$ into $H'$.

\subsection{Separated metric observations and atomic refinement}

The reciprocal scalar geometry supports an explicit test of its clock
dictionary. On the computed weak geometry preparation, all seven
ground-to-excited ratios of the complete eight-level local Hamiltonians
differ from one another and from the metric lapse ratio. The devices
have fixed coordinate extent and changing proper dimensions, so this
does not constitute an equivalence-principle refutation. It establishes
that this device choice supplies no universal metric clock attachment.
A declared transverse scalar-wave extension uses the same density
$C(r)$, retaining all 31 normal modes and the vacuum. Its complete
detector laws have total-variation distances from the flat construction
of approximately $4.0$ to $7.7$ parts per million at the executed times. The associated
piecewise metric has transverse ray delay
$6.4197332251\,10^{-5}$ in source time and radial null speed one.
The dispersive wave arrival is not identified with the ray. A harmonic
scalar tail also has nonzero linearized Einstein residual under the
tested vacuum or pressureless interpretation; the regional Einstein
theorem therefore cannot supply that interpretation without its
stress and geometry premises.

For the separate Dirichlet atom with supplied $m=1$ and $g^2=4\pi$,
volume and spacing effects are tested independently. At fixed spacing
$a=1/3$, boxes of side $24$, $32$ and $40$ give first gaps
$0.420953171076$, $0.420395862457$ and $0.420371341048$.
At fixed side $24$, spacings $1/3$, $1/4$, $3/16$ and $3/20$
give $0.420953171076$, $0.399138718909$, $0.388151976665$ and
$0.383407843476$, compared with the hydrogenic reference $3/8$.
The second-shell splitting decreases from $0.005500452345$ to
$0.000410646363$. Exact reflection reduction represents even standing
waves over all $159^3=4{,}019{,}679$ sites through $512{,}000$
even-sector rows; the other parity sectors are separate.
The discrete maximum principle, zero extension and min--max theorem
prove monotonicity of each ordered parity eigenvalue with volume at
fixed spacing. They do not prove monotonicity of gaps or spacing
convergence. Residuals bound distance to some spectrum; the reported
floating residuals and successive-grid differences are not certified
continuum uncertainties or level-index proofs.
